
\documentclass[10pt]{article}
\usepackage[utf8]{inputenc}
\usepackage[T1]{fontenc}
\usepackage{amsmath, amssymb, xcolor, geometry, enumitem, multicol, tcolorbox}

\definecolor{mainblue}{RGB}{0, 102, 204}
\geometry{a4paper, margin=1.2cm, top=1cm, bottom=3cm, footskip=1.2cm}

\begin{document}
\enlargethispage{2.5cm}

% Modern Header
\begin{tcolorbox}[colback=mainblue!5, colframe=mainblue, sharp corners, boxrule=0.5pt]
    \small \textbf{Unit:} Unit 0: Foundations of Algebra \hfill \textbf{Name:} \rule{3cm}{0.4pt} \\
    \small \textbf{Topic:} Variables and Expressions / Order of Operations \hfill \textbf{Date:} \rule{2cm}{0.4pt} \\
    \centering \textbf{\large Homework (Jalapeno)}
\end{tcolorbox}

\vspace{0.2cm}

% MCQ Section
\begin{multicols}{2}
\begin{enumerate}[label=\textbf{Q\arabic*.}, leftmargin=*, itemsep=6pt, parsep=0pt]

  \item \small Evaluate the expression: $2 \cdot 3 + 12 - 8$
  \begin{enumerate}[label=(\Alph*), leftmargin=0.6cm, nosep, itemsep=2pt]
    \item 10
    \item 12
    \item 14
    \item 16
    \item 18
  \end{enumerate}

  \item \small Translate the phrase into an algebraic expression: '5 more than the product of 3 and x'
  \begin{enumerate}[label=(\Alph*), leftmargin=0.6cm, nosep, itemsep=2pt]
    \item $3x + 5$
    \item $3x - 5$
    \item $5x + 3$
    \item $5x - 3$
    \item $3 + 5x$
  \end{enumerate}

  \item \small Simplify the expression: $2x + 5x - 3x$
  \begin{enumerate}[label=(\Alph*), leftmargin=0.6cm, nosep, itemsep=2pt]
    \item $4x$
    \item $6x$
    \item $8x$
    \item $10x$
    \item $12x$
  \end{enumerate}

  \item \small Substitute x = 4 into the expression: $x^2 + 2x - 3$
  \begin{enumerate}[label=(\Alph*), leftmargin=0.6cm, nosep, itemsep=2pt]
    \item 13
    \item 15
    \item 17
    \item 19
    \item 21
  \end{enumerate}

  \item \small Evaluate the expression: $12 - 3 + 2 \cdot 4$
  \begin{enumerate}[label=(\Alph*), leftmargin=0.6cm, nosep, itemsep=2pt]
    \item 13
    \item 15
    \item 17
    \item 19
    \item 21
  \end{enumerate}

  \item \small The ancient Egyptians used the expression $5x - 2$ to calculate the area of a rectangular field. If x = 6, what is the area?
  \begin{enumerate}[label=(\Alph*), leftmargin=0.6cm, nosep, itemsep=2pt]
    \item 28
    \item 30
    \item 32
    \item 34
    \item 36
  \end{enumerate}

  \item \small The Babylonians used the expression $2x + 5$ to calculate the perimeter of a triangular field. If x = 3, what is the perimeter?
  \begin{enumerate}[label=(\Alph*), leftmargin=0.6cm, nosep, itemsep=2pt]
    \item 11
    \item 13
    \item 15
    \item 17
    \item 19
  \end{enumerate}

  \item \small The ancient Greeks used the expression $x^2 + 2x - 6$ to calculate the volume of a cube. If x = 2, what is the volume?
  \begin{enumerate}[label=(\Alph*), leftmargin=0.6cm, nosep, itemsep=2pt]
    \item 0
    \item 2
    \item 4
    \item 6
    \item 8
  \end{enumerate}

\end{enumerate}
\end{multicols}

\vspace{-0.2cm}
\hrule
\vspace{0.2cm}
\textbf{\small Open Ended Questions (FRQ)}
\begin{enumerate}[label=\textbf{Q\arabic*.}, start=9, itemsep=5pt, topsep=0pt]

  \item \small Evaluate the expression: $3 \cdot (2x + 1) - 2(x - 3)$, given x = 2 \vspace{2cm}

  \item \small Simplify the expression: $(x + 2)(x - 3)$, and then substitute x = 4 \vspace{2cm}

\end{enumerate}

\vfill

% Chef's Tips Footer
\begin{tcolorbox}[colback=blue!5, colframe=mainblue!50, title=\small \textbf{Chef's Tips}, fonttitle=\bfseries, bottom=1mm]
\begin{itemize}[noitemsep, topsep=2pt, leftmargin=0.5cm]
  \item \scriptsize Remind students to use the correct order of operations\item \scriptsize Encourage students to read the questions carefully\item \scriptsize Allow students to use a calculator to check their answers
\end{itemize}
\end{tcolorbox}

\end{document}
  